\section{What the evidence changes about memory design}
\label{sec:discussion}
\subsection{A useful abstraction with distinct failure modes}
SKILL.state separates working context from the transcript. The coding
experiments support that architectural idea while showing why a smaller
state is not a complete memory policy. Three measured effects matter.
Sparse runtime records leave useful execution evidence absent; retaining
commands and changes improves on file references alone. Removing recent
interactions increases the pilot's repetition diagnostic and lowers
success in the larger comparisons. Changing state and history also changes
cache reuse, so logical token reduction is not a price comparison.

The remedies address different parts of this problem. Runtime ownership
avoids asking the model to rewrite every observed fact. A recent window
preserves detail that a schema did not select. Their combination still
requires evaluation: enriching the runtime record did not produce
monotone gains, and history recovery had different cost consequences in
the two implementations. The evidence supports a division of memory
responsibilities, not a claim that the tested hybrid dominates full history.

The extended window sweep turns $k=6$ from an unexplained constant into
a measured threshold: success and repeated invocation stop improving there,
but at the ceiling and floor of a 12-task pilot. Wider windows keep adding
prompt without removing the per-turn cost of rewriting state. Some lost
information therefore concerns the local execution sequence; older
evidence may still matter on tasks that leave headroom, and a
late-relevance task could need an observation outside every finite window.

\subsection{Why fixed-horizon savings do not settle this question}
The original synthetic scaling experiments hold the event stream fixed;
our agent chooses actions and when to stop~\citep{skillstate2026}.
A context bound at a fixed horizon does not constrain the extra work
a memory policy may induce. The original paper also reports public
interactive benchmarks, so this distinction is not an explanation of all
differences between the studies. Different models, schemas, and runtimes
remain alternative explanations. Its reported total-token savings likewise
cannot determine monetary savings without input, output, and cache
accounting. Supplementary~\ref{app:pricing} illustrates this missing
information without claiming to reconstruct the authors' bill.

The horizon cohorts provide no universal rule to switch from transcript
to state. State alone loses success in every cohort, whereas limited
history recovers much of the gap even in the long-reference cohort.
An adaptive policy should therefore respond to information needs and
complete-task cost, not assume that trajectory age makes recent evidence
dispensable. A retrievable archive and a matched history-only control are
the next tests: the former checks recovery of older evidence, the latter
whether state adds value beyond a recent window.

\paragraph{Scope.}
These are implementation-specific findings on one model and task
collection. Exploratory configuration selection, in-sample schema
adaptation, unadjusted repeated-task tests, incomplete execution-source
provenance, and aggregate-only hybrid reporting limit inference.
The pilot supports descriptive resource and repetition comparisons, not a
success-equivalence test. Neither causal forgetting nor the effectiveness
of adaptive switching and retrieval is established.
